\chapter{Complexity — Relativization}
%%% generated by the blueprint pipeline from TCSlib.Complexity.Diagonalization.Relativization
\section{Overview}

Declarations of the Lean module \texttt{TCSlib.Complexity.Diagonalization.Relativization}.

\section{Declarations}

\begin{definition}[Unary witness language]
\lean{Complexity.unaryWitnessLang}
\statementsource{BGS75}{pdf-2459e3840724-p006-b001}
\label{Complexity.unaryWitnessLang}
\leanok
Given a language $B$ of Boolean strings, define a language consisting of those Boolean
strings $w$ for which there is a natural number $n$ such that $w$ is the length-$n$
string of true bits and there is a string $x$ of length $n$ belonging to $B$.
\end{definition}

\begin{definition}[Nondeterministic oracle table]
\lean{Complexity.uwTable}
\statementsource{BGS75}{
  pdf-2459e3840724-p002-b001,
  pdf-2459e3840724-p002-b002
}
\label{Complexity.uwTable}
\leanok
\uses{Turing.OracleNDTM, Turing.OracleTM.answerAction}
Define a nondeterministic oracle Turing machine with no ordinary work tapes, Boolean
symbols, and four states numbered $0,1,2,3$. Its initial state is $0$, its query state
is $1$, and its positive and negative answer states are $2$ and $3$. In state $0$, on
input symbol $1$, it nondeterministically writes its choice bit to the query tape, moves
right, and remains in state $0$; on input symbol $0$, it takes the negative-answer
action. On a blank input cell it moves right and enters the query state. In state $2$ it
takes the positive-answer action, and in every other state it takes the negative-answer
action.
\end{definition}

\begin{definition}[Finite nondeterministic oracle machine]
\lean{Complexity.uwTM}
\statementsource{BGS75}{
  pdf-2459e3840724-p002-b001,
  pdf-2459e3840724-p002-b002
}
\label{Complexity.uwTM}
\leanok
\uses{Complexity.uwTable, Turing.FinOracleNDTM}
Define a finite nondeterministic oracle Turing machine over the Boolean alphabet with no
ordinary work tapes and four states numbered $0,1,2,3$. Its underlying machine is the
table specified by the initial state $0$, query state $1$, positive-answer state $2$,
and negative-answer state $3$; these three special states are pairwise distinct.
\end{definition}

\begin{definition}[Initial scanning configuration]
\lean{Complexity.uwScanCfg}
\label{Complexity.uwScanCfg}
\leanok
\uses{Turing.Cfg, Turing.FinTM.bufferTape}
For a Boolean input string $x$, an index $j$ with $j\le |x|$, and a Boolean string $u$,
define a configuration with state $0$, input-head position $j+1$ in the shifted indexing
of the input, an empty output string, and its sole work tape initialized from $u$ by the
buffer-tape construction, with the work-head position at $j$.
\end{definition}

\begin{definition}[Initial query configuration]
\lean{Complexity.uwQueryCfg}
\label{Complexity.uwQueryCfg}
\leanok
\uses{Turing.Cfg, Turing.FinTM.bufferTape}
For a Boolean input string $x$ and a Boolean string $u$, define a configuration with
state $1$, input-head position $|x|+1$ in the shifted indexing of the input, an empty
output string, and its sole work tape initialized from $u$ by the buffer-tape
construction, with the work-head position at $|x|$.
\end{definition}

\begin{theorem}[Initial scan configuration]
\lean{Complexity.uw_init}
\label{Complexity.uw_init}
\leanok
\uses{Complexity.uwScanCfg, Complexity.uwTable, Turing.OracleNDTM.initCfg}
\difficulty{3}
For every Boolean string $x$, the machine's initial configuration on input $x$ is the
scan configuration at position $0$ with an empty buffer.
\end{theorem}

\begin{theorem}[Ordinary transition outside the query state]
\lean{Complexity.uw_step_scan}
\label{Complexity.uw_step_scan}
\leanok
\uses{Complexity.uwTable, Turing.Action.apply, Turing.Cfg, Turing.Cfg.inputSymbol, Turing.Cfg.workTapeSymbols, Turing.OracleNDTM.stepWith}
\difficulty{3}
Let $B$ be a language of bit strings, and let $c$ be a configuration of the designated
oracle machine whose state is its state $0$. For either choice bit $b$, one step with
oracle $B$ is exactly the configuration obtained by applying the machine's transition
action for choice $b$, state $0$, and the symbols currently scanned on the input and
work tapes.
\end{theorem}

\begin{theorem}[One step of the unary scan]
\lean{Complexity.uw_scan_step}
\label{Complexity.uw_scan_step}
\leanok
\uses{Complexity.uwScanCfg, Complexity.uwTable, Complexity.uw_step_scan, Turing.Action.apply, Turing.FinTM.bufferTape_append, Turing.OracleNDTM.stepWith, Turing.inputSymbolInner, Turing.moveInputPos_pos_of_ne_right}
\difficulty{4}
Let $B$ be a language of bit strings, let $x$ be a bit string, and let $j<|x|$ with the
bit of $x$ at position $j$ equal to $1$. If $u$ is a bit string of length $j$ and $b$ is
a bit, then the designated oracle machine, started in its scan configuration at position
$j$ with buffer $u$, takes one step with choice $b$ to the scan configuration at
position $j+1$ with buffer $u$ followed by $b$.
\end{theorem}

\begin{theorem}[A one-choice run is one step]
\lean{Complexity.uw_runWith_single}
\label{Complexity.uw_runWith_single}
\leanok
\uses{Complexity.uwTable, Turing.Cfg, Turing.OracleNDTM.runWith, Turing.OracleNDTM.stepWith}
\difficulty{1}
Let $B$ be a language of Boolean strings, let $b$ be a Boolean choice, let $x$ be a
Boolean input string, and let $c$ be any machine configuration on input $x$. Running the
machine from $c$ with the one-bit choice string $[b]$ gives the same configuration as
taking one step from $c$ using choice $b$.
\end{theorem}

\begin{theorem}[Run through a true input prefix]
\lean{Complexity.uw_scan_run}
\label{Complexity.uw_scan_run}
\leanok
\uses{Complexity.uwScanCfg, Complexity.uwTable, Complexity.uw_init, Complexity.uw_runWith_single, Complexity.uw_scan_step, Turing.OracleNDTM.initCfg, Turing.OracleNDTM.runWith, Turing.OracleNDTM.runWith_append, Turing.OracleNDTM.runWith_nil}
\difficulty{4}
Let $B$ be a language of Boolean strings, and let $x,w$ be Boolean strings. For any
$j\le |x|$ such that every bit of $x$ at an index less than $j$ is true, if $j\le |w|$,
then running the machine on the prefix of $w$ of length $j$ from its initial
configuration on input $x$ yields the scan configuration at position $j$, with buffer
equal to that same prefix.
\end{theorem}

\begin{theorem}[Transition from scanning to querying]
\lean{Complexity.uw_end_step}
\label{Complexity.uw_end_step}
\leanok
\uses{Complexity.uwQueryCfg, Complexity.uwScanCfg, Complexity.uwTable, Complexity.uw_step_scan, Turing.Action.apply, Turing.Cfg.inputSymbol, Turing.OracleNDTM.stepWith}
\difficulty{3}
Let $B$ be a language of Boolean strings, and let $x,u$ be Boolean strings and $b$ a
Boolean choice. One machine step from the scan configuration at position $|x|$, with
buffer $u$, yields the query configuration with buffer $u$.
\end{theorem}

\begin{theorem}[Rejecting at a false input bit]
\lean{Complexity.uw_false_step}
\label{Complexity.uw_false_step}
\leanok
\uses{Complexity.uwScanCfg, Complexity.uwTable, Complexity.uw_step_scan, Turing.OracleNDTM.stepWith, Turing.inputSymbolInner}
\difficulty{3}
Let $B$ be a language of Boolean strings, let $x,u$ be Boolean strings, and let $j$ be
an index with $j<|x|$ and $x_j=0$. For either Boolean choice $b$, one machine step from
the scan configuration at position $j$ with buffer $u$ halts and outputs the one-bit
string $[0]$.
\end{theorem}

\begin{theorem}[The query routine's output]
\lean{Complexity.uw_tail}
\label{Complexity.uw_tail}
\leanok
\uses{Complexity.uwQueryCfg, Complexity.uwTable, Turing.OracleNDTM.fixChoice, Turing.OracleNDTM.runWith, Turing.OracleNDTM.runWith_cons, Turing.OracleNDTM.runWith_nil, Turing.OracleNDTM.stepWith_eq_fixChoice_step, Turing.OracleTM.queryString, Turing.OracleTM.queryString_bufferTape, Turing.OracleTM.step_of_answerAction, Turing.OracleTM.step_of_qQuery}
\difficulty{4}
Let $B$ be a language of bit strings, let $x$ and $u$ be bit strings, and let $b_1,b_2$
be bits. Starting from the designated query configuration with buffer $u$ on input $x$,
the machine run using choices $b_1,b_2$ halts and outputs the single bit $1$ if $u\in
B$, and $0$ otherwise.
\end{theorem}

\begin{theorem}[Acceptance on a fully true input]
\lean{Complexity.uw_branch}
\label{Complexity.uw_branch}
\leanok
\uses{Complexity.uwQueryCfg, Complexity.uwScanCfg, Complexity.uwTable, Complexity.uw_end_step, Complexity.uw_false_step, Complexity.uw_runWith_single, Complexity.uw_scan_run, Complexity.uw_tail, Turing.OracleNDTM.initCfg, Turing.OracleNDTM.runWith, Turing.OracleNDTM.runWith_append, Turing.OracleNDTM.runWith_of_halt, Turing.OracleNDTM.stepWith}
\difficulty{6}
Let $B$ be a language of Boolean strings, and let $x,w$ be Boolean strings with
$|x|+3\le |w|$. After the machine runs on the choice string $w$ from its initial
configuration on input $x$, it is halted; its output is the one-bit string $[1]$ exactly
when every bit of $x$ is true and the prefix of $w$ of length $|x|$ belongs to $B$.
\end{theorem}

\begin{theorem}[Unary witnesses are in nondeterministic polynomial time]
\lean{Complexity.unaryWitnessLang_mem_NPOracle}
\label{Complexity.unaryWitnessLang_mem_NPOracle}
\leanok
\uses{Complexity.NPOracle, Complexity.unaryWitnessLang, Complexity.uwTM, Complexity.uw_branch}
\difficulty{5}
Let $B$ be a language of finite Boolean words. The language consisting of the words that
are all true bits of some length $n$, for which there exists a word $x$ of length $n$
belonging to $B$, is in nondeterministic polynomial time with oracle $B$.
\end{theorem}

\begin{definition}[Relabeling an oracle machine's states]
\lean{Complexity.oracleRelabel}
\label{Complexity.oracleRelabel}
\leanok
\uses{Turing.Action.mapState, Turing.OracleTM}
Given an oracle Turing machine with state set $S$ and an equivalence $e:S\simeq S'$,
define the machine over state set $S'$ by applying $e$ to its initial, query,
positive-answer, and negative-answer states. Its transition at state $q'$, scanned input
symbol, and scanned work-tape symbols is the original transition at $e^{-1}(q')$, with
each successor state relabeled by $e$.
\end{definition}

\begin{theorem}[Relabeling preserves well-formedness]
\lean{Complexity.oracleRelabel_wellFormed}
\label{Complexity.oracleRelabel_wellFormed}
\leanok
\uses{Complexity.oracleRelabel, Turing.OracleTM, Turing.OracleTM.WellFormed}
\difficulty{1}
Let $M$ be an oracle Turing machine with state set $S$, and let $e:S\to S'$ be a
bijection. If the query, positive-answer, and negative-answer states of $M$ are pairwise
distinct, then the machine obtained by relabeling every designated state and every
transition successor via $e$ also has pairwise distinct query, positive-answer, and
negative-answer states.
\end{theorem}

\begin{theorem}[State relabeling commutes with one oracle step]
\lean{Complexity.oracleRelabel_step}
\label{Complexity.oracleRelabel_step}
\leanok
\uses{Complexity.oracleRelabel, Turing.Cfg, Turing.Cfg.inputSymbol, Turing.Cfg.mapState, Turing.Cfg.mapState_apply, Turing.Cfg.workTapeSymbols, Turing.OracleTM, Turing.OracleTM.queryString, Turing.OracleTM.step}
\difficulty{5}
Let $k\in\mathbb{N}$, let $\Sigma$ be an alphabet, let $S,S'$ be state sets, let $x$ be
a word over $\Sigma$, let $M$ be an oracle Turing machine with $k$ ordinary work tapes
and states $S$, and let $e\colon S\simeq S'$ be a bijection. For every oracle language
$O$ over $\Sigma$ and configuration $c$ on input $x$, one step of the relabeled machine
from the relabeling of $c$ equals the relabeling of the configuration reached by one
step of $M$ from $c$.
\end{theorem}

\begin{theorem}[State relabeling commutes with oracle runs]
\lean{Complexity.oracleRelabel_runFrom}
\label{Complexity.oracleRelabel_runFrom}
\leanok
\uses{Complexity.oracleRelabel, Complexity.oracleRelabel_step, Turing.Cfg, Turing.Cfg.mapState, Turing.OracleTM, Turing.OracleTM.runFrom, Turing.OracleTM.step}
\difficulty{3}
Let $k\in\mathbb{N}$, let $\Sigma$ be an alphabet, let $S,S'$ be state sets, let $x$ be
a word over $\Sigma$, let $M$ be an oracle Turing machine with $k$ ordinary work tapes
and states $S$, and let $e\colon S\simeq S'$ be a bijection. For every oracle language
$O$ over $\Sigma$, configuration $c$ on input $x$, and $t\in\mathbb{N}$, running the
relabeled machine for $t$ steps from the configuration obtained by relabeling $c$ gives
the configuration obtained by relabeling the result of running $M$ from $c$ for $t$
steps.
\end{theorem}

\begin{theorem}[State relabeling preserves timed computation]
\lean{Complexity.oracleRelabel_computesInTime}
\label{Complexity.oracleRelabel_computesInTime}
\leanok
\uses{Complexity.oracleRelabel, Complexity.oracleRelabel_runFrom, Turing.Cfg.mapState, Turing.OracleTM, Turing.OracleTM.ComputesInTime, Turing.OracleTM.initCfg}
\difficulty{3}
Let $k$ be a natural number, let $\Sigma$ be an alphabet, let $S$ and $S'$ be state
sets, let $M$ be an oracle Turing machine with $k$ ordinary work tapes, and let $e\colon
S\simeq S'$ be a bijection. For every oracle language $O$ over $\Sigma$, input word $x$,
output word $y$, and time bound $t\in\mathbb{N}$, the machine obtained by relabeling the
states of $M$ via $e$ computes $y$ from $x$ relative to $O$ within $t$ steps if and only
if $M$ does.
\end{theorem}

\begin{definition}[Canonical well-formed oracle tables]
\lean{Complexity.CanonTab}
\label{Complexity.CanonTab}
\leanok
\uses{Turing.OracleTM, Turing.OracleTM.WellFormed}
For natural numbers $k$ and $n$, define the type of canonical tables to consist of
Boolean oracle Turing machines with $k$ ordinary work tapes and state set
$\{0,\ldots,n-1\}$, equipped with a proof that their query, positive-answer, and
negative-answer states are pairwise distinct.
\end{definition}

\begin{definition}[Finite Type of Canonical Oracle Machines]
\lean{Complexity.canonTabFintype}
\label{Complexity.canonTabFintype}
\leanok
\uses{Complexity.CanonTab, Turing.OracleTM}
For natural numbers $k$ and $n$, the collection of well-formed oracle Turing machines
with $k$ ordinary work tapes, Boolean tape symbols, and states indexed by $\text{Fin}(n)$
is finite.
\end{definition}

\begin{definition}[Default finite oracle machine]
\lean{Complexity.oracleDefault}
\label{Complexity.oracleDefault}
\leanok
\uses{Turing.FinOracleTM}
Define a finite oracle Turing machine over the Boolean alphabet with no ordinary work
tapes and three states numbered $0,1,2$. Its initial state and positive-answer state are
both $1$, its query state is $0$, and its negative-answer state is $2$. At every state
and on every scanned input and work-tape symbol, its transition leaves the heads
stationary, makes no writes or output, and halts; the query, positive-answer, and
negative-answer states are pairwise distinct.
\end{definition}

\begin{definition}[Oracle table at a given index]
\lean{Complexity.oracleTableAt}
\label{Complexity.oracleTableAt}
\leanok
\uses{Complexity.CanonTab, Complexity.canonTabFintype, Complexity.oracleDefault, Turing.FinOracleTM}
For natural numbers $k$, $n$, and $r$, let $C_{k,n}$ be the finite type of canonical
oracle-machine tables. If $r<|C_{k,n}|$, the definition returns the well-formed Boolean
oracle Turing machine represented by the table at index $r$ in the enumeration of
$C_{k,n}$; if $r\ge |C_{k,n}|$, it returns the default Boolean oracle Turing machine.
\end{definition}

\begin{definition}[Decoded oracle table machine]
\lean{Complexity.decodeOracleTable}
\label{Complexity.decodeOracleTable}
\leanok
\uses{Complexity.oracleTableAt, Turing.FinOracleTM}
For each natural number $j$, decode the oracle machine obtained by first unpairing $j$,
then unpairing the second coordinate, and using the resulting three natural numbers as
the work-tape count, state count, and table index, respectively.
\end{definition}

\begin{theorem}[Ranking and Decoding Oracle Machines]
\lean{Complexity.decodeOracleTable_rank}
\label{Complexity.decodeOracleTable_rank}
\leanok
\uses{Complexity.CanonTab, Complexity.canonTabFintype, Complexity.decodeOracleTable, Complexity.oracleTableAt, Turing.FinOracleTM}
\difficulty{3}
Let $k,n$ be natural numbers and let a well-formed oracle Turing machine with $k$
ordinary work tapes, Boolean alphabet, and state set of size $n$ be given. If this
machine is assigned its index in the finite enumeration of such machines, then decoding
the natural-number code obtained by pairing $k$, $n$, and that index returns the same
finite oracle machine, with its three distinguished states and transition function
unchanged.
\end{theorem}

\begin{definition}[Oracle machine enumeration]
\lean{Complexity.oracleEnumeration}
\statementsource{BGS75}{
  pdf-2459e3840724-p003-b006,
  pdf-2459e3840724-p004-b004
}
\label{Complexity.oracleEnumeration}
\leanok
\uses{Complexity.decodeOracleTable, Turing.FinOracleTM}
For each natural number $i$, take the first coordinate in the fixed pairing
decomposition of $i$ and decode that natural number as an oracle table machine.
\end{definition}

\begin{theorem}[Pairing preserves the decoded oracle table]
\lean{Complexity.oracleEnumeration_pair}
\label{Complexity.oracleEnumeration_pair}
\leanok
\uses{Complexity.decodeOracleTable, Complexity.oracleEnumeration}
\difficulty{1}
For all natural numbers $j$ and $r$, the oracle machine with index obtained by pairing
$j$ and $r$ is the machine decoded from $j$.
\end{theorem}

\begin{theorem}[Enumeration of Finite Oracle Machines]
\lean{Complexity.exists_finOracleTM_enumeration}
\label{Complexity.exists_finOracleTM_enumeration}
\leanok
\uses{Complexity.CanonTab, Complexity.canonTabFintype, Complexity.decodeOracleTable, Complexity.decodeOracleTable_rank, Complexity.oracleEnumeration, Complexity.oracleEnumeration_pair, Complexity.oracleRelabel, Complexity.oracleRelabel_computesInTime, Complexity.oracleRelabel_wellFormed, Turing.FinOracleTM, Turing.FinOracleTM.ComputesInTime}
\difficulty{5}
There is a sequence $(N_i)_{i\in\mathbb N}$ of finite oracle Turing machines over the
binary alphabet such that, for every finite binary oracle machine $M$ and every natural
number $i_0$, there is an index $i\ge i_0$ for which $N_i$ and $M$ compute exactly the
same outputs within exactly $t$ steps, for every binary oracle language, binary input
and output strings, and natural number $t$.
\end{theorem}

\begin{theorem}[A Missing Bit String of Prescribed Length]
\lean{Complexity.exists_length_not_mem}
\label{Complexity.exists_length_not_mem}
\leanok
\difficulty{6}
Let $n$ be a natural number and let $Q$ be a finite list of binary strings with
$|Q|<2^n$. Then there is a binary string $z$ of length $n$ that does not occur in $Q$.
\end{theorem}

\begin{theorem}[A Length Exceeding Exponential and Polynomial Bounds]
\lean{Complexity.exists_fresh_len}
\label{Complexity.exists_fresh_len}
\leanok
\difficulty{5}
For every pair of natural numbers $i,b$, there is a natural number $n$ such that $b<n$,
$i+2\le n$, and $n^i+i<2^n$.
\end{theorem}

\begin{definition}[Stage data]
\lean{Complexity.StageData}
\label{Complexity.StageData}
\leanok
A stage datum consists of a finite set of finite Boolean strings and a natural-number
bound.
\end{definition}

\begin{definition}[Positive finite oracle]
\lean{Complexity.posOracle}
\statementsource{BGS75}{
  pdf-2459e3840724-p006-b003,
  pdf-2459e3840724-p006-b004
}
\label{Complexity.posOracle}
\leanok
For a finite set $P$ of finite Boolean strings, the associated language consists exactly
of the strings in $P$.
\end{definition}

\begin{definition}[Fresh length]
\lean{Complexity.freshLen}
\statementsource{BGS75}{
  pdf-2459e3840724-p006-b003,
  pdf-2459e3840724-p006-b004
}
\label{Complexity.freshLen}
\leanok
\uses{Complexity.exists_fresh_len}
For natural numbers $i$ and $b$, choose a natural number witnessing the property that
there is a fresh length for $i$ and $b$.
\end{definition}

\begin{theorem}[Properties of the chosen length]
\lean{Complexity.freshLen_spec}
\label{Complexity.freshLen_spec}
\leanok
\uses{Complexity.exists_fresh_len, Complexity.freshLen}
\difficulty{1}
For natural numbers $i$ and $b$, let $m$ be the length chosen so that $b<m$, $i+2\le m$,
and $m^i+i<2^m$. Then these three inequalities hold for $m$.
\end{theorem}

\begin{definition}[Stage query list]
\lean{Complexity.stageQ}
\statementsource{BGS75}{
  pdf-2459e3840724-p006-b003,
  pdf-2459e3840724-p006-b004
}
\label{Complexity.stageQ}
\leanok
\uses{Complexity.StageData, Complexity.freshLen, Complexity.posOracle, Turing.FinOracleTM, Turing.OracleTM.queriesWithin}
Let $N$ be a sequence of finite Boolean oracle machines, let $i$ be a natural number,
and let $d$ be stage data with finite set $P$ and bound $b$. The stage query list is the
list of queries made by machine $N(i)$, using the oracle consisting of $P$, on the
all-true Boolean input of length $L$, within $L^i+i$ steps, where $L$ is the fresh
length chosen for $i$ and $b$.
\end{definition}

\begin{theorem}[Bound on the stage query list]
\lean{Complexity.stageQ_length}
\label{Complexity.stageQ_length}
\leanok
\uses{Complexity.StageData, Complexity.freshLen, Complexity.freshLen_spec, Complexity.stageQ, Turing.FinOracleTM, Turing.OracleTM.queriesWithin_length_le}
\difficulty{1}
Let $N$ be a sequence of finite oracle Turing machines, let $i$ be a natural number, and
let $d$ be stage data. The number of queries in the list produced at stage $i$ using $d$
is strictly less than $2^{\text{freshLen}(i,d.\text{bound})}$.
\end{theorem}

\begin{definition}[Stage witness string]
\lean{Complexity.stageWit}
\statementsource{BGS75}{
  pdf-2459e3840724-p006-b003,
  pdf-2459e3840724-p006-b004
}
\label{Complexity.stageWit}
\leanok
\uses{Complexity.StageData, Complexity.exists_length_not_mem, Complexity.stageQ_length, Turing.FinOracleTM}
Let $N$ be a sequence of finite Boolean oracle machines, let $i$ be a natural number,
and let $d$ be stage data. Choose a Boolean string whose length is not the length of any
query in the stage query list for $N$, $i$, and $d$.
\end{definition}

\begin{theorem}[Length and avoidance of the stage witness]
\lean{Complexity.stageWit_spec}
\label{Complexity.stageWit_spec}
\leanok
\uses{Complexity.StageData, Complexity.exists_length_not_mem, Complexity.freshLen, Complexity.stageQ, Complexity.stageQ_length, Complexity.stageWit, Turing.FinOracleTM}
\difficulty{1}
Let $N$ be a sequence of finite oracle Turing machines, let $i$ be a natural number, and
let $d$ be stage data. The witness selected for this stage is a Boolean string of length
$\text{freshLen}(i,d.\text{bound})$ that does not occur among the queries made by the
$i$th machine during this stage.
\end{theorem}

\begin{definition}[Acceptance at a stage]
\lean{Complexity.StageAccepts}
\statementsource{BGS75}{
  pdf-2459e3840724-p006-b003,
  pdf-2459e3840724-p006-b004
}
\label{Complexity.StageAccepts}
\leanok
\uses{Complexity.StageData, Complexity.freshLen, Complexity.posOracle, Turing.FinOracleTM, Turing.FinOracleTM.ComputesInTime}
For a sequence $N$ of finite oracle Turing machines with Boolean input symbols, a stage
index $i$, and stage data $d$, acceptance means that the $i$th machine, with oracle
consisting of the positive strings in $d$, halts on the all-true input of length
$\text{freshLen}(i,d.\text{bound})$ within $\text{freshLen}(i,d.\text{bound})^i+i$ steps and
outputs the one-symbol string $[\mathrm{true}]$.
\end{definition}

\begin{definition}[One stage update]
\lean{Complexity.stageStep}
\statementsource{BGS75}{
  pdf-2459e3840724-p006-b003,
  pdf-2459e3840724-p006-b004
}
\label{Complexity.stageStep}
\leanok
\uses{Complexity.StageAccepts, Complexity.StageData, Complexity.freshLen, Complexity.stageWit, Turing.FinOracleTM}
Given a sequence $N$ of finite oracle Turing machines over Boolean symbols, an index
$i\in\mathbb N$, and stage data $d$, the updated stage data has positive set equal to
the current positive set if the $i$th machine accepts at stage $i$; otherwise it adjoins
the selected witness string. Its bound is the maximum of
$\text{freshLen}(i,d.\text{bound})$ and $\text{freshLen}(i,d.\text{bound})^i+i$.
\end{definition}

\begin{definition}[Iterated stage construction]
\lean{Complexity.stage}
\statementsource{BGS75}{
  pdf-2459e3840724-p006-b003,
  pdf-2459e3840724-p006-b004
}
\label{Complexity.stage}
\leanok
\uses{Complexity.StageData, Complexity.stageStep, Turing.FinOracleTM}
For a sequence $N$ of finite oracle Turing machines over Boolean symbols, define stage
data recursively by starting with empty positive set and bound $0$ at stage $0$, then
obtaining stage $i+1$ by applying the stage update at index $i$ to stage $i$.
\end{definition}

\begin{definition}[Stage-generated oracle language]
\lean{Complexity.bgsOracle}
\statementsource{BGS75}{
  pdf-2459e3840724-p006-b003,
  pdf-2459e3840724-p006-b004
}
\label{Complexity.bgsOracle}
\leanok
\uses{Complexity.stage, Turing.FinOracleTM}
For a sequence $N$ of finite oracle Turing machines over Boolean symbols, the associated
oracle language consists of exactly those Boolean strings that belong to the positive
set at some stage of the recursive stage construction.
\end{definition}

\begin{theorem}[Successive Positive Sets Are Nested]
\lean{Complexity.stage_pos_succ}
\label{Complexity.stage_pos_succ}
\leanok
\uses{Complexity.stage, Complexity.stageStep, Turing.FinOracleTM}
\difficulty{4}
Let $(N_j)_{j\ge0}$ be a sequence of finite-state oracle Turing machines, and let $i$ be
a nonnegative integer. Every string positive at stage $i$ is positive at stage $i+1$.
\end{theorem}

\begin{theorem}[Monotonicity of Positive Sets]
\lean{Complexity.stage_pos_mono}
\label{Complexity.stage_pos_mono}
\leanok
\uses{Complexity.stage, Complexity.stage_pos_succ, Turing.FinOracleTM}
\difficulty{4}
Let $(N_j)_{j\ge0}$ be a sequence of finite-state oracle Turing machines, and let $i\le
j$ be nonnegative integers. Every string positive at stage $i$ is positive at stage $j$.
\end{theorem}

\begin{theorem}[Successor stage bound]
\lean{Complexity.stage_bound_succ}
\label{Complexity.stage_bound_succ}
\leanok
\uses{Complexity.freshLen, Complexity.stage, Turing.FinOracleTM}
\difficulty{1}
For every sequence $N$ of finite oracle Turing machines and every $i\in\mathbb{N}$, the
bound at stage $i+1$ is the maximum of the fresh length chosen from $i$ and the bound at
stage $i$, and the $i$th power of that stage bound plus $i$.
\end{theorem}

\begin{theorem}[Monotonicity of stage bounds]
\lean{Complexity.stage_bound_mono}
\label{Complexity.stage_bound_mono}
\leanok
\uses{Complexity.freshLen_spec, Complexity.stage, Complexity.stage_bound_succ, Turing.FinOracleTM}
\difficulty{4}
For every sequence $N$ of finite oracle Turing machines and natural numbers $i\le j$,
the bound at stage $i$ is at most the bound at stage $j$.
\end{theorem}

\begin{theorem}[Growth of Later Stage Lengths]
\lean{Complexity.stage_later}
\label{Complexity.stage_later}
\leanok
\uses{Complexity.freshLen, Complexity.freshLen_spec, Complexity.stage, Complexity.stage_bound_mono, Complexity.stage_bound_succ, Turing.FinOracleTM}
\difficulty{3}
Let $(N_j)_{j\ge0}$ be a sequence of finite-state oracle Turing machines, and let $i<j$
be nonnegative integers. If $n_r$ denotes the fresh length chosen at stage $r$ using
that stage's bound, then $n_i^i+i<n_j$ and $n_i<n_j$.
\end{theorem}

\begin{theorem}[New Positive Strings at a Stage]
\lean{Complexity.stage_new}
\label{Complexity.stage_new}
\leanok
\uses{Complexity.StageAccepts, Complexity.stage, Complexity.stageStep, Complexity.stageWit, Turing.FinOracleTM}
\difficulty{4}
Let $(N_j)_{j\ge0}$ be a sequence of finite-state oracle Turing machines, let $i$ be a
nonnegative integer, and let $z$ be a Boolean string. If $z$ is positive at stage $i+1$
but not positive at stage $i$, then $z$ is the witness selected at stage $i$, and
machine $N_i$ does not accept its specified unary input within $n_i^i+i$ steps relative
to the stage-$i$ positive-set oracle, where $n_i$ is the fresh length for that stage.
\end{theorem}

\begin{theorem}[First entry into the stage sets]
\lean{Complexity.exists_entry}
\label{Complexity.exists_entry}
\leanok
\uses{Complexity.bgsOracle, Complexity.stage, Turing.FinOracleTM}
\difficulty{4}
Let $N=(N_i)_{i\in\mathbb N}$ be a sequence of finite oracle Turing machines over the
binary alphabet, with each machine’s query, positive-answer, and negative-answer states
pairwise distinct. Let $P_i$ be the positive set at stage $i$ of the stage construction
determined by $N$, and let $B=\bigcup_{i\in\mathbb N}P_i$. Every binary word $z\in B$
belongs to $P_{s+1}$ but not to $P_s$ for some $s\in\mathbb N$.
\end{theorem}

\begin{theorem}[Length of a Newly Added Word]
\lean{Complexity.entry_length}
\label{Complexity.entry_length}
\leanok
\uses{Complexity.freshLen, Complexity.stage, Complexity.stageWit_spec, Complexity.stage_new, Turing.FinOracleTM}
\difficulty{3}
Let $N$ be a sequence of finite oracle Turing machines over the binary alphabet, and let
$s$ be a natural number. If a binary word $z$ belongs to the positive set at stage $s+1$
but not to the positive set at stage $s$, then its length is the fresh length determined
by $s$ and the bound stored in the stage data at stage $s$.
\end{theorem}

\begin{theorem}[Consistency of Stage Queries]
\lean{Complexity.stage_consistent}
\label{Complexity.stage_consistent}
\leanok
\uses{Complexity.bgsOracle, Complexity.entry_length, Complexity.exists_entry, Complexity.posOracle, Complexity.stage, Complexity.stageQ, Complexity.stageWit_spec, Complexity.stage_later, Complexity.stage_new, Complexity.stage_pos_mono, Turing.FinOracleTM, Turing.OracleTM.length_le_of_mem_queriesWithin}
\difficulty{5}
Let $(N_i)_{i\ge 0}$ be a sequence of finite-state oracle Turing machines, and let $S_i$
be the finite set of positive strings at stage $i$. Use the stage-$i$ oracle whose
language is $S_i$, and run $N_i$ on the unary input of length $n_i$, where $n_i$ is the
fresh length chosen for stage $i$ and its current bound; let $Q_i$ be the list of query
strings it makes during the first $n_i^i+i$ steps. Every $z\in Q_i$ belongs to $S_i$ if
and only if it belongs to $S_j$ for some stage $j$.
\end{theorem}

\begin{theorem}[Transfer Between Stage and Limit Oracles]
\lean{Complexity.stage_transfer}
\label{Complexity.stage_transfer}
\leanok
\uses{Complexity.bgsOracle, Complexity.freshLen, Complexity.posOracle, Complexity.stage, Complexity.stage_consistent, Turing.FinOracleTM, Turing.FinOracleTM.ComputesInTime, Turing.OracleTM.ComputesInTime, Turing.OracleTM.runFrom_eq_of_agree_queriesWithin}
\difficulty{3}
Let $(N_j)_{j\ge0}$ be a sequence of finite-state oracle Turing machines, let $i$ be a
nonnegative integer, and let $y$ be a Boolean string. Running $N_i$ with the union of
all stage-positive sets as oracle on the unary input of length $n_i$ halts within
$n_i^i+i$ steps with output $y$ if and only if running it with the stage-$i$
positive-set oracle on that same input halts within the same number of steps with output
$y$, where $n_i$ is the fresh length chosen for stage $i$.
\end{theorem}

\begin{theorem}[Unary witness strings of fixed length]
\lean{Complexity.mem_unaryWitnessLang_replicate}
\label{Complexity.mem_unaryWitnessLang_replicate}
\leanok
\uses{Complexity.unaryWitnessLang}
\difficulty{4}
Let $B$ be a language of bit strings and let $n\in\mathbb{N}$. The unary string
consisting of $n$ copies of $1$ belongs to the unary witness language of $B$ if and only
if $B$ contains a bit string of length $n$.
\end{theorem}

\begin{theorem}[Acceptance and Unary Witness Exclusion]
\lean{Complexity.stage_flip}
\label{Complexity.stage_flip}
\leanok
\uses{Complexity.StageAccepts, Complexity.bgsOracle, Complexity.entry_length, Complexity.exists_entry, Complexity.freshLen, Complexity.mem_unaryWitnessLang_replicate, Complexity.stage, Complexity.stageStep, Complexity.stageWit, Complexity.stageWit_spec, Complexity.stage_later, Complexity.stage_new, Complexity.stage_transfer, Complexity.unaryWitnessLang, Turing.FinOracleTM, Turing.FinOracleTM.ComputesInTime}
\difficulty{5}
Let $(N_j)_{j\ge0}$ be a sequence of finite-state oracle Turing machines and let $i$ be
a nonnegative integer. Write $n_i$ for the fresh length chosen at stage $i$ from its
stage bound, and let $B$ be the language of strings that are unary words of some length
$n$ for which a length-$n$ Boolean string belongs to the oracle. Then $N_i$, using the
union of all stage-positive sets as its oracle, accepts the unary input of length $n_i$
within $n_i^i+i$ steps and outputs the one-symbol word $[\mathrm{true}]$ if and only if
the unary input of length $n_i$ does not belong to $B$.
\end{theorem}

\begin{theorem}[Stage Time Bound Dominates]
\lean{Complexity.stage_dominates}
\label{Complexity.stage_dominates}
\leanok
\uses{Complexity.freshLen, Complexity.freshLen_spec, Complexity.stage, Turing.FinOracleTM}
\difficulty{4}
Let $(N_j)_{j\ge0}$ be a sequence of finite-state oracle Turing machines, and let
$c,k,i$ be nonnegative integers satisfying $k+1+2c\le i$. If $n_i$ is the fresh length
chosen at stage $i$ using that stage's bound, then $c(n_i^k+1)\le n_i^i+i$.
\end{theorem}

\begin{theorem}[An oracle separating nondeterministic and deterministic polynomial time]
\lean{Complexity.exists_oracle_ne}
\label{Complexity.exists_oracle_ne}
\leanok
\uses{Complexity.NPOracle, Complexity.POracle, Complexity.bgsOracle, Complexity.exists_finOracleTM_enumeration, Complexity.freshLen, Complexity.stage, Complexity.stage_dominates, Complexity.stage_flip, Complexity.unaryWitnessLang, Complexity.unaryWitnessLang_mem_NPOracle, Turing.FinOracleTM.ComputesInTime, Turing.FinOracleTM.ComputesInTime.mono, Turing.FinOracleTM.ComputesInTime.output_unique, Turing.MultiTapeTM.indicator}
\difficulty{6}
There exists a language $B$ of binary strings such that the language consisting of
strings $1^n$ for which some binary string of length $n$ belongs to $B$ is in
nondeterministic polynomial time relative to oracle $B$, but is not in deterministic
polynomial time relative to oracle $B$.
\end{theorem}

\begin{theorem}[Oracle separation and equality]
\lean{Complexity.baker_gill_solovay}
\label{Complexity.baker_gill_solovay}
\leanok
\uses{Complexity.EXPCOM, Complexity.NPOracle, Complexity.POracle, Complexity.POracle_EXPCOM_eq_NPOracle_EXPCOM, Complexity.exists_oracle_ne}
\difficulty{4}
There is a Boolean language $A$ for which deterministic and nondeterministic polynomial
time with oracle $A$ coincide, and there is a Boolean language $B$ for which these two
oracle complexity classes differ.
\end{theorem}
